\section{Related Work}
\label{sec:related}

\paragraph{Partial rotation and frequency selection.}
Slow RoPE bands can support content matching rather than distinct position
responses \citep{barbero2025round}. MHA2MLA uses partial rotation and frequency
selection as part of converting pretrained attention to a compressed latent
architecture \citep{ji2025mha2mla}. Thus neither reducing rotary dimensions
nor the importance of selection at a small width is a new observation here.
Our proposed comparison holds the complete attention architecture fixed and
asks whether an anchored analytic allocation improves on explicit high-frequency
and uniform-subset baselines under the same training budget.

\paragraph{Learned and data-dependent frequency use.}
LeRoPE learns frequencies and compares them with geometric and partial RoPE
\citep{karypis2026lerope}; frequency usage also depends on the data's dependency
structure \citep{wu2026datashapes}. These results rule out treating a static
spectral diagnostic as a universal utility metric. Our candidate uses one fixed
analytic shape without a frequency search, but that distinction alone does not
establish a practical advantage. The six-arm experiment must determine it.

\paragraph{Context extension.}
Position interpolation, YaRN, LongRoPE, and MrRoPE modify the treatment of
positions or frequencies during extension
\citep{chen2024position,peng2024yarn,ding2024longrope,tian2026mrrope}.
MrRoPE studies mixed-radix transports of an existing table; the present
experiment studies training-time allocation under changing rotary width.
These mathematical design spaces can overlap. MrRoPE is relevant extension
prior work rather than a substitute for the budget controls here. Any future
zero-training superiority claim would require its own matched strong-baseline
comparison.
